\section{Fitted Q-Iteration 与深度 Q 网络}

\subsection{从表格到函数近似}

当状态空间连续或很大时，无法维护 Q 值表格。解决方案是用函数近似器（如神经网络）来参数化 Q 函数。

\begin{importantbox}{Fitted Q-Iteration (FQI)}
FQI 将 Q-Learning 转化为监督学习问题：
\begin{enumerate}
    \item 用当前 Q 网络计算 TD target：$y_i = r_i + \gamma \max_{a'} Q_{\theta^-}(s_i', a')$
    \item 将 $(s_i, a_i)$ 作为输入，$y_i$ 作为标签，训练 Q 网络：
    $$\min_\theta \sum_i \left(Q_\theta(s_i, a_i) - y_i\right)^2$$
    \item 更新目标网络 $\theta^- \leftarrow \theta$
    \item 重复
\end{enumerate}
\end{importantbox}

\subsection{DQN 的关键技巧}

\begin{importantbox}{DQN (Deep Q-Network) 的创新}
DQN 引入了两个关键技巧使深度 Q-Learning 稳定工作：
\begin{enumerate}
    \item \textbf{Experience Replay}：将历史经验存储在 replay buffer 中，训练时随机采样 mini-batch。这打破了数据的时间相关性，提高了样本利用率。
    \item \textbf{Target Network}：使用一个周期性更新的目标网络 $Q_{\theta^-}$ 来计算 TD target。这减少了"移动的目标"问题带来的不稳定性。
\end{enumerate}
\end{importantbox}

\begin{warningbox}{函数近似 Q-Learning 的不稳定性}
与表格型不同，使用函数近似的 Q-Learning \textbf{不保证收敛}。Bellman 更新的收缩性质在函数近似下可能失效，导致：
\begin{itemize}
    \item Q 值发散（overestimation bias 的累积）
    \item 训练不稳定（oscillation 或 divergence）
    \item 对超参数敏感
\end{itemize}
这就是为什么需要 target network、experience replay 等稳定化技巧。
\end{warningbox}

\subsection{DQN 的改进}

\begin{itemize}
    \item \textbf{Double DQN}：用当前网络选择动作，目标网络评估值，减少 overestimation
    \item \textbf{Dueling DQN}：将 Q 函数分解为 $V(s) + A(s, a)$，分别估计状态值和优势
    \item \textbf{Prioritized Experience Replay}：按 TD error 大小优先采样，关注更有学习价值的经验
\end{itemize}

\subsection{本章小结}

从表格 Q-Learning 到 DQN 的核心跨越是函数近似。Experience replay 和 target network 是确保训练稳定的关键技巧。

